\begin{proof}
Para simplificar, pongamos $A'=\widehat A$ y consideremos $t$ como un
elemento de $A'$; se tiene $A'/tA'=(A/tA)\otimes_AA'$, y, como $A'$ es un
$A$-módulo plano, de (3.3.1) se deduce que los ideales primos de $A'$ asociados
al $A'$-módulo $A'/tA'$ son los ideales primos $\mathfrak q'_{jh}$ tales que
$\mathfrak q'_{jh}$ está por encima de $\mathfrak q_j$ y está asociado al
$k_j$-módulo $(A'/tA')\otimes_Ak_j$, donde $k_j$ es el cuerpo de fracciones
de $A/\mathfrak q_j$. Ahora bien,
$\operatorname{Spec}((A'/tA')\otimes_Ak_j)$ es la fibra formal de $A$ en el
punto $\mathfrak q_j$, o también la de $A/\mathfrak q_j$ en el punto
$\mathfrak q_j$ (punto genérico de
$\operatorname{Spec}(A/\mathfrak q_j)$). Puesto que, en virtud de (iv), el
completado $A'/\mathfrak q_jA'$ de $A/\mathfrak q_j$ es reducido, lo mismo
ocurre con la fibra formal de $A/\mathfrak q_j$ en el punto genérico de
$\operatorname{Spec}(A/\mathfrak q_j)$; por tanto, esta fibra formal no tiene
ningún ciclo primo asociado inmerso y sus anillos locales en los puntos
genéricos $\mathfrak q'_{jh}$ de sus componentes irreducibles son cuerpos
(3.2.1). Así pues, por (3.3.3) y la hipótesis (ii), $A'/tA'$ no tiene ideales
primos asociados inmersos y, por otra parte,
$\mathfrak q_jA'_{\mathfrak q'_{jh}}$ es el ideal maximal de
$A'_{\mathfrak q'_{jh}}$ para todo $h$; como $A_{\mathfrak q_j}$ es un anillo
de valoración discreta, su ideal maximal $\mathfrak q_jA_{\mathfrak q_j}$
es principal, y, por tanto, el ideal maximal del anillo local noetheriano
$A'_{\mathfrak q'_{jh}}$ es principal. Esto implica que dicho anillo es un
anillo de valoración discreta (Bourbaki, \emph{Alg. comm.}, cap. VI, § 3,
n\textsuperscript{o} 6, prop. 9). Acabamos, pues, de verificar las hipótesis
(i) a (iv) para \emph{el anillo local completo} $A'$.
Basta, por tanto, demostrar que, si $A$ es \emph{completo} y satisface las hipótesis
(i) a (iii) (pues (iv) se satisface automáticamente en este caso), entonces
$A$ es reducido. Ahora bien, las hipótesis (i) y (ii) implican que, si
$\varphi:A\longrightarrow\prod_{j=1}^nA_{\mathfrak q_j}$
es el homomorfismo canónico, se tiene
$t^nA=\varphi^{-1}(\varphi(t^nA))$ (3.4.9); como la imagen canónica de $tA$
en el anillo de valoración discreta $A_{\mathfrak q_j}$ es una potencia
$\mathfrak q_j^hA_{\mathfrak q_j}$ del ideal maximal, se ve que, en $A$, la
topología $(tA)$-preádica es la imagen inversa por $\varphi$ de la
topología producto en $\prod_{j=1}^nA_{\mathfrak q_j}$. Ahora bien, como $t$
pertenece al radical de Jacobson de $A$, la topología $(tA)$-preádica es
separada; por tanto, $\varphi$ es \emph{inyectivo}. Pero, en virtud de la hipótesis
(iii), los $A_{\mathfrak q_j}$ son dominios de integridad, de modo que
$\prod_{j=1}^nA_{\mathfrak q_j}$ es reducido y \emph{a fortiori} también lo
es $A$. Queda demostrado.
\end{proof}
\end{proof}
